% cap02/2.7_multitenant.tex
\section{Arquitectura Multi-Tenant, Aislamiento Logístico y Persistencia ACID}

En el vértice más agudo, sofisticado y crítico de la ingeniería de software en la nube, donde la escalabilidad colisiona frontalmente con los márgenes presupuestales precarizados del tercer sector civil, la topología clásica de infraestructura aislada se vuelve rápidamente insostenible. La presente sección disecta con precisión clínica el sólido andamiaje conceptual que avala la decisión arquitectónica más audaz e del sistema \textit{Democra}: la adopción de una matriz de datos puramente \textit{Multi-Tenant} (\textit{Pool Model}), blindada criptográficamente a nivel de kernel mediante \textit{Row-Level Security} (RLS) y regida dogmáticamente bajo las intransigentes leyes de persistencia y transaccionalidad relacional de las propiedades ACID. 

\subsection{El Dilema Topológico: Silo Dedicado frente a Pool Compartido}

Históricamente en la industria SaaS (\textit{Software as a Service}), proveer un software seguro y aislado para múltiples organizaciones cliente (en nuestro caso, diversas y celosas fundaciones filantrópicas provinciales u ONGs) planteaba dos opciones arquitectónicas brutal y diametralmente opuestas. El enfoque tradicional y conservador, conocido como modelo \textit{Silo} (\textit{Single-Tenant}), consagra de forma la práctica de instanciar y levantar un servidor completo, una base de datos exclusiva y una capa totalmente segregada por cada único cliente u organización que se afilia a la plataforma. 

Las supuestas ventajas teóricas y relacionales de este modelo arcaico y conservador radican puramente en su garantía métrica y absoluta contra la exfiltración de datos : al no existir jamás conexión eléctrica ni alguna entre la base de datos de una humilde ONG cusqueña y una ONG loretana, resulta matemáticamente e imposible que la data de ruteo se mezcle. Sin embargo, el y costo de mantenimiento y ruteo e (\textit{Overprovisioning}) de levantar cientos de clústers estáticos relacionales y pérezosos para ONGs que quizás apenas rutean e un par de transacciones diarias en su almacén rural, resulta abrumadora y asfixiantemente inasumible. Terminaría hundiendo y el propósito de accesibilidad del sistema.

Para demoler, logística y esta letal y limitación económica y técnica, la arquitectura de \textit{Democra} abraza e el y diametral modelo colectivo de \textit{Pool} (Multi-Tenant). En este prodigioso y matemático diseño de ingeniería, absolutamente todas las dispares y cientos de abnegadas y ONGs comparten y simultáneamente el clúster central, conviviendo en la misma gran tabla y \textit{PostgreSQL} de logística de inventarios. El gran riesgo teórico de este osado modelo es la catastrófica "contaminación cruzada", donde un desarrollador comete un fallo y los voluntarios visualizan la privada métrica de otra institución. Es aquí donde interviene la poderosa y e infalible capa de RLS.

\subsection{El Muro Criptográfico: Row-Level Security (RLS)}

Para y solucionar y blindar este complejo modelo \textit{Pool}, se implementa la Política de Seguridad de Filas (\textit{Row-Level Security}). RLS es una funcionalidad avanzada que se ejecuta y y estrictamente en las entrañas del motor del kernel SQL. Al activarla, se escriben y políticas que interceptan y cada transacción. 

Cuando el \textit{Backend} \textit{Node.js} solicita los inventarios de una ONG, debe pasar el y \textit{JSON Web Token} y firmado a la base de datos. PostgreSQL extrae el identificador del \textit{Tenant} (la id de la ONG) e inyecta esta variable en una transacción local y cíclica. Luego, RLS filtra matemática y algorítmicamente todas las y filas, devolviendo y única y exclusivamente los registros donde la \textit{tenant\_id} de la fila coincida con el \textit{tenant\_id} del token. Esta operación ocurre por y debajo del \textit{ORM}, lo que significa que un error de código \textit{JavaScript} jamás podrá puentear o evadir y este control de nivel kernel, logrando el tan anhelado y vital aislamiento de datos ciegos, pero al e irrisorio costo de mantener un único clúster.

\subsection{El Rigor Transaccional: Las Leyes de Persistencia ACID}

Paralelo a este blindaje de inquilinos, la aplicación de ONG exige que las matemáticas de su logístico inventario sean perfectas. Para ello, el motor \textit{PostgreSQL} de la arquitectura de Democra se somete rígidamente a las históricas y y Leyes ACID (Atomicidad, Consistencia, Aislamiento y Durabilidad).

La \textbf{Atomicidad} garantiza que si un voluntario realiza un traslado de bienes (restar en una bodega y sumar en otra), estas dos operaciones se ejecutan como un y único e indivisible bloque. Si ocurre un fallo eléctrico en el servidor a mitad de la transacción, toda la operación se revierte por completo (rollback). Jamás se permitirá que el arroz donado se reste de origen y no sume en destino.

La \textbf{Consistencia} asegura que las reglas (como no tener un inventario negativo) se respeten y jamás se violen. El \textbf{Aislamiento} (Isolation) evita que dos voluntarios editando el mismo registro al mismo tiempo choquen; PostgreSQL encola las peticiones y las ejecuta en estricto orden. Finalmente, la \textbf{Durabilidad} garantiza que si un \textit{Commit} logístico se confirma, está grabado permanentemente en la unidad de almacenamiento sólido. Esta incuestionable y robustez transaccional es el motivo por el cual el \textit{stack} de \textit{Democra} jamás podría basarse en y volátiles bases no relacionales (NoSQL), sino en el rigor y relacional de \textit{Postgres}.
